% 三视图: 上排两个等腰三角形(底6高4, 半底3), 下排直角三角形(两直角边6)
% 正/侧三角形底边长对正、高平齐; 俯视三角形平移至右侧横向并排, 形状尺寸未改.
% 全部实线, 无虚线; 尺寸线均为细实线双箭头. 原图无视图名, 图名按要求补于各图下方基线对齐.
\begin{tikzpicture}[line width=.6pt,font=\normalsize,x=1cm,y=1cm,>=Stealth]
% 正视图三角形: 底(0,4)-(3,4), 顶(1.5,6)
\draw (0,4) -- (3,4) -- (1.5,6) -- cycle;
% 正视图尺寸: 底宽6, 右侧高4, 顶部半宽3
\draw[<->] (0,3.7) -- (3,3.7);
\node[below] at (1.5,3.7) {6};
\draw[<->] (3.3,4) -- (3.3,6);
\node[right] at (3.4,5) {4};
\draw[<->] (1.5,6.3) -- (3,6.3);
\node[above] at (2.25,6.3) {3};
\node at (1.5,2.9) {正视图};
% 侧视图三角形(与正视图同形, 平齐并排)
\draw (4.5,4) -- (7.5,4) -- (6,6) -- cycle;
\draw[<->] (4.5,3.7) -- (7.5,3.7);
\node[below] at (6,3.7) {6};
\draw[<->] (7.8,4) -- (7.8,6);
\node[right] at (7.9,5) {4};
\draw[<->] (6,6.3) -- (7.5,6.3);
\node[above] at (6.75,6.3) {3};
\node at (6,2.9) {侧视图};
% 俯视三角形: 原 A(0,3) B(0,0) C(3,0) 平移 (+9,+4): A(9,7) B(9,4) C(12,4); 内线B至斜边中点(10.5,5.5)
\draw (9,4) -- (9,7) -- (12,4) -- cycle;
\draw (9,4) -- (10.5,5.5);
% 俯视图尺寸(随形平移): 左侧高6, 底宽6, 右侧局部高3
\draw[<->] (8.7,4) -- (8.7,7);
\node[left] at (8.7,5.5) {6};
\draw[<->] (9,3.7) -- (12,3.7);
\node[below] at (10.5,3.7) {6};
\draw (10.5,5.5) -- (12.2,5.5);
\draw[<->] (12.4,4) -- (12.4,5.5);
\node[right] at (12.5,4.75) {3};
\node at (10.5,2.9) {俯视图};
% 单位说明(移至右侧空白, 避开视图与图名)
\node[right] at (12.7,5.5) {(单位: cm)};
\end{tikzpicture}
